\label{triptych:brief-synthesis:start}
\section*{The Propers: Themes and Movement}

\subsection*{Hunger disclosed}

Israel's complaint is not romanticized: hunger and fear are real, yet memory
has contracted Egypt to full pots and forgotten slavery. The crowd in John
also seeks Jesus after being fed. Chrysostom reads Christ's correction as a
diagnosis of motive, not a contempt for bodily need. The liturgy begins with
people gathered and strengthened by God, then permits the readings to expose
how a gift can be received without yet recognizing its giver.

\subsection*{Gift before comprehension}

Manna arrives before Israel can name it. Psalm 78 retells that gift so another
generation will remember. John intensifies the movement: the Father, not
Moses, gives the true bread, and the gift is not merely an object but the Son
who gives life to the world. Augustine distinguishes the sign that fills the
stomach from the reality to which faith comes. The Gospel's final
self-identification makes the Lectionary correlation explicit without
collapsing manna and Christ into identical things.

\subsection*{Learning Christ}

Ephesians is semi-continuous, not selected as a historical commentary on
John. Heard together, however, it guards against a merely acquisitive account
of faith. To believe the one sent is to have learned Christ: the old person is
put off, the mind renewed, and the new person put on. Chrysostom's homily on
Ephesians emphasizes that renewal reaches the governing faculty of life and
appears in truthful conduct. The acclamation likewise refuses to isolate
physical bread from the word proceeding from God.

\subsection*{Received bread, renewed life}

The Eucharistic units keep agency in the order of grace. Gifts already
received from divine bounty are returned; the sacrament is received as the
perpetual memorial of the Passion; its saving fruit is requested rather than
presumed. Communion B repeats John~6:35 and therefore closes the most direct
bread sequence. Communion A instead calls the soul to remember all God's
benefits. Either enacted path joins reception to gratitude and renewal, but
the guide does not pretend the two optional antiphons were simultaneously
used.

\clearpage

The movement can therefore be scanned in four stages:

\begin{fourstable}{Stage}{Textual movement}{Reception control}{Limit}
Need & Murmuring and searching disclose bodily hunger and disordered memory. &
Chrysostom, \emph{Homilies on John} 43--44 & Bodily need is not despised.\\
Provision & Manna is given, remembered, and redirected toward the Father's
true bread. & Augustine on Psalm 77 and John~6 & Type and fulfillment are not
made identical.\\
Renewal & Belief in the sent Son entails putting off the old person and
putting on the new. & Chrysostom, \emph{Homily on Ephesians} 13 & The apostolic
reading remains semi-continuous.\\
Communion & Gift returns as offering and is received as Passion memorial with
fruit sought. & Ambrose, \emph{On the Mysteries} 8 & Sacramental fruit is not
treated as automatic.\\
\end{fourstable}

The source-grounded center is neither ``earn the bread'' nor ``ignore earthly
hunger.'' Exodus shows gratuitous provision accompanied by a daily discipline;
John distinguishes the perishing from the enduring while naming the Son as
giver and gift; Ephesians describes the kind of life shaped by learning him.
The Collect's contrast between passing and enduring goods therefore concerns
their right use, not rejection of creation. The Communion alternatives then
lead either through remembered benefits or directly through Christ's bread
saying.

\begin{studybox}{Relationship control}
The First Reading and Gospel are officially correlated. The Psalm is
responsorial, the acclamation acclamatory, and Ephesians belongs to the
semi-continuous apostolic course. Connections to the Week XVIII prayers and
antiphons are textual observations and source-grounded synthesis, not claims
about historical selection intent.
\end{studybox}
\label{triptych:brief-synthesis:end}
